\documentclass[11pt]{article}
\usepackage{mathblatt}


\begin{document}
\pruefheftstil{Prozent}{Prüfungsheft P10 \textperiodcentered{} schwach}
\noindent{\Large\bfseries Prozent}\hfill{\small\color{mbgrau}Prüfungsheft P10 \textperiodcentered{} schwach}\par\medskip
\pfrueckblick
\begin{pfnr}{}{1.}{}
In einer Tüte sind $10$ Bonbons. $6$ davon sind rot. Schreibe den Anteil der roten Bonbons als Bruch und kürze ihn.
\par\noindent \leerfeld \par
\rechenplatz{3}
\fusshilfe{1: $\tfrac{3}{5}$}
\end{pfnr}
\begin{pfnr}{}{2.}{}
Erweitere den Bruch $\tfrac{7}{10}$ auf den Nenner $100$.
\par\noindent \leerfeld \par
\fusshilfe{2: $\tfrac{70}{100}$}
\end{pfnr}
\begin{pfnr}{}{3.}{}
Schreibe $\tfrac{1}{2}$ als Dezimalzahl.
\par\noindent \leerfeld \par
\fusshilfe{3: $0{,}5$}
\end{pfnr}
\begin{pfnr}{}{4.}{}
Schreibe $0{,}9$ als Bruch.
\par\noindent \leerfeld \par
\rechenplatz{2}
\fusshilfe{4: $\tfrac{9}{10}$}
\end{pfnr}
\begin{pfnr}{}{5.}{}
Berechne. $\tfrac{600}{100}$
\par\noindent \leerfeld \par
\fusshilfe{5: $6$}
\end{pfnr}
\begin{pfnr}{}{6.}{}
Berechne. $0{,}1 \cdot 70$
\par\noindent \leerfeld \par
\rechenplatz{2}
\fusshilfe{6: $7$}
\end{pfnr}
\begin{pfnr}{}{7.}{}
$2$ Kinokarten kosten $18\,€$. Berechne mit der Tabelle, wie viel $5$ Kinokarten kosten.
\par\smallskip\noindent \begin{dreisatz}{Karten}{Preis}\dsz{2}{18\,€}\dsp{\phantom{:0}}{\phantom{:0}}\dsz{1}{\dsleer}\dsp{\phantom{:0}}{\phantom{:0}}\dsz{5}{\dsleer}\end{dreisatz}\par
\par\noindent \leerfeld[€] \par
\rechenplatz{4}
\fusshilfe{7: $45\,€$}
\end{pfnr}
\begin{pfnr}{}{8.}{}
$4$ kg Äpfel kosten $12\,€$. Berechne mit der Tabelle, wie viel $3$ kg Äpfel kosten.
\par\smallskip\noindent \begin{dreisatz}{Gewicht}{Preis}\dsz{4\,\text{kg}}{12\,€}\dsp{\phantom{:0}}{\phantom{:0}}\dsz{1\,\text{kg}}{\dsleer}\dsp{\phantom{:0}}{\phantom{:0}}\dsz{3\,\text{kg}}{\dsleer}\end{dreisatz}\par
\par\noindent \leerfeld[€] \par
\rechenplatz{4}
\fusshilfe{8: $9\,€$}
\end{pfnr}
\begin{pfnr}{}{9.}{}
Runde $4{,}73$ auf eine Stelle nach dem Komma.
\par\noindent \leerfeld \par
\rechenplatz{2}
\fusshilfe{9: $4{,}7$}
\end{pfnr}
\begin{pfnr}{}{10.}{}
Runde $8{,}26$ auf eine Stelle nach dem Komma.
\par\noindent \leerfeld \par
\rechenplatz{2}
\fusshilfe{10: $8{,}3$}
\end{pfnr}
\begin{pfnr}{}{11.}{}
Berechne $\tfrac{1}{4}$ von $20\,€$.
\par\noindent \leerfeld[€] \par
\fusshilfe{11: $5\,€$}
\end{pfnr}
\begin{pfnr}{}{12.}{}
Berechne $\tfrac{1}{2}$ von $36$ kg.
\par\noindent \leerfeld[kg] \par
\fusshilfe{12: $18$ kg}
\end{pfnr}
\begin{pfnr}{}{13.}{}
Berechne. $2^{6}$
\par\noindent \leerfeld \par
\fusshilfe{13: 64}
\end{pfnr}
\begin{pfnr}{}{14.}{}
Berechne. $5^{3}$
\par\noindent \leerfeld \par
\fusshilfe{14: 125}
\end{pfnr}
\pfstufekopf[sprozentundanteilumwa]{Prozent und Anteil umwandeln}{in 1 der letzten 5 Prüfungen}
\pfgruppe{}
\begin{pfnr}{eigene Aufgabe}{15.}{}
In einer Stadt fahren $20\,\%$ der Berufstätigen mit dem Bus zur Arbeit. Ergänze: Jeder \leerfeld. Berufstätige fährt mit dem Bus.
\rechenplatz{3}
\fusshilfe{15: jeder $5.$}
\end{pfnr}
\pfgruppe{}
\begin{pfnr}{eigene Aufgabe}{16.}{}
Schreibe den Bruch $\tfrac{11}{20}$ in Prozent.
\par\noindent \leerfeld[\%] \par
\rechenplatz{3}
\fusshilfe{16: $55\,\%$}
\end{pfnr}
\begin{pfnr}{eigene Aufgabe}{17.}{}
Der Streifen zeigt ein Ganzes. Ein Teil davon ist grau. Wie viel Prozent des ganzen Streifens sind grau?
\par\smallskip\noindent \streifenfeld{50}{$0\,\%$}{$100\,\%$}\par
\par\noindent \leerfeld[\%] \par
\rechenplatz{2}
\fusshilfe{17: $50\,\%$}
\end{pfnr}
\begin{pfnr}{eigene Aufgabe}{18.}{}
Der Streifen zeigt ein Ganzes. Ein Teil davon ist grau. Wie viel Prozent des ganzen Streifens sind grau?
\par\smallskip\noindent \streifenfeld{20}{$0\,\%$}{$100\,\%$}\par
\par\noindent \leerfeld[\%] \par
\rechenplatz{2}
\fusshilfe{18: $20\,\%$}
\end{pfnr}
\pfgruppe{}
\begin{pfnr}{P10 ’22}{19.}{1 BE}
Ein Fünftel aller Jugendlichen erhält kein Taschengeld. Kreuze den passenden Prozentsatz an:
\pfkreuzzeile{5 \%}\pfkreuzzeile{20 \%}\pfkreuzzeile{25 \%}\pfkreuzzeile{50 \%}
\fusshilfe{19: 20 \%}
\end{pfnr}
\begin{pfnr}{P10 ’20}{20.}{1 BE}
Ein Förster sagt: „Bei mir wachsen nur 4 \% der gepflanzten Bäume nicht an.“ Kreuze an, welche Aussage dasselbe bedeutet:
\pfkreuzzeile{Von 10 Bäumen wachsen 4 nicht an.}\pfkreuzzeile{Jeder vierte Baum wächst nicht an.}\pfkreuzzeile{Von 100 Bäumen wachsen 4 nicht an}
\fusshilfe{20: 4 von 100 Bäumen wachsen nicht an}
\end{pfnr}
\pfstufekopf[sprozentwert]{Prozentwert $W$}{in 1 der letzten 5 Prüfungen}
\pfgruppe{}
\begin{pfnr}{eigene Aufgabe}{21.}{}
Teile den Streifen in gleiche Teile ein. Jeder Teil soll $10\,\%$ sein.
\par\smallskip\noindent \streifenleer[0]\par
\fusshilfe{21: $10$ gleiche Teile}
\end{pfnr}
\begin{pfnr}{eigene Aufgabe}{22.}{}
Teile den Streifen in gleiche Teile ein. Jeder Teil soll $25\,\%$ sein.
\par\smallskip\noindent \streifenleer[0]\par
\fusshilfe{22: $4$ gleiche Teile}
\end{pfnr}
\begin{pfnr}{eigene Aufgabe}{23.}{}
In der Klasse sind $28$ Kinder. $7$ davon haben einen Hund. Welche Zahl ist hier das Ganze?
\fusshilfe{23: $28$ Kinder}
\end{pfnr}
\begin{pfnr}{eigene Aufgabe}{24.}{}
Von $600$ Losen sind $90$ Gewinne. Welche Zahl ist hier das Ganze?
\fusshilfe{24: $600$ Lose}
\end{pfnr}
\pfgruppe{}
\begin{pfnr}{eigene Aufgabe}{25.}{}
Der ganze Streifen steht für $48\,€$. Wie viel Euro sind $25\,\%$ davon?
\par\smallskip\noindent \streifen[4]{25}{$0$}{$48\,€$}\par
\par\noindent \leerfeld[€] \par
\rechenplatz{3}
\fusshilfe{25: $12\,€$}
\end{pfnr}
\begin{pfnr}{eigene Aufgabe}{26.}{}
Der ganze Streifen steht für $64\,€$. Wie viel Euro sind $25\,\%$ davon?
\par\smallskip\noindent \streifen[4]{25}{$0$}{$64\,€$}\par
\par\noindent \leerfeld[€] \par
\rechenplatz{3}
\fusshilfe{26: $16\,€$}
\end{pfnr}
\begin{pfnr}{eigene Aufgabe}{27.}{}
Eine Tischlampe kostet $64\,€$. Es gibt $25\,\%$ Rabatt. Wie viel kostet die Lampe jetzt?
\par\noindent \leerfeld[€] \par
\rechenplatz{4}
\fusshilfe{27: $48\,€$}
\end{pfnr}
\pfab{$W = G \cdot p$}
\begin{pfnr}{eigene Aufgabe}{28.}{}
Berechne $20\,\%$ von $4{,}50\,€$.
\par\noindent \leerfeld[€] \par
\rechenplatz{3}
\fusshilfe{28: $0{,}90\,€$}
\end{pfnr}
\begin{pfnr}{eigene Aufgabe}{29.}{}
Bücher kosten ohne Mehrwertsteuer $40\,€$. Dazu kommen $7\,\%$ Mehrwertsteuer. Wie viel Euro Mehrwertsteuer sind das?
\par\noindent \leerfeld[€] \par
\rechenplatz{3}
\fusshilfe{29: $2{,}80\,€$}
\end{pfnr}
\begin{pfnr}{P10 ’26}{30.}{1 BE}
Gib 30 \% von 70 € an.
\rechenplatz{2}
\fusshilfe{30: 21 €, Tipp: $W$ = 70 € · 0,3}
\end{pfnr}
\begin{pfnr}{P10 ’14}{31.}{1 BE}
Bestimme 13 \% von 50 €.
\rechenplatz{2}
\fusshilfe{31: 6,50 €}
\end{pfnr}
\pfgruppe{}
\begin{pfnr}{eigene Aufgabe}{32.}{}
Dieselbe Jacke kostet in Laden A $95\,€$ mit $20\,\%$ Rabatt. In Laden B kostet sie $86\,€$ mit $10\,\%$ Rabatt. In welchem Laden ist die Jacke billiger?
\rechenplatz{5}
\fusshilfe{32: $1{,}40\,€$, in Laden A ist die Jacke billiger}
\end{pfnr}
\begin{pfnr}{P10 ’19}{33.}{2 BE}
In einer Studie wurden dieselben 1 200 Jugendlichen über mehrere Jahre befragt. 2013 hatten 72 \% von ihnen ein Smartphone, 2016 waren es 95 \%. Berechne, um wie viele Jugendliche die Zahl der Smartphone-Besitzer in dieser Zeit gestiegen ist.
\pfzwfrage{Gib an, um wie viele Prozentpunkte der Anteil der Smartphone-Besitzer gestiegen ist.}
\rechenplatz{4}
\fusshilfe{33: 276 Jugendliche}
\end{pfnr}
\begin{pfnr}{P10 ’17}{34.}{1 BE}
Für eine Bohrmaschine zum Preis von 120,00 € gibt es an der Kasse 20 \% Nachlass. Gib an, wie viel Euro der Nachlass beträgt.
\rechenplatz{2}
\fusshilfe{34: 24 €}
\end{pfnr}
\pfgruppe{}
\begin{pfnr}{P10 ’21}{35.}{1 BE}
Ein Fahrrad kostet 550 €. Wer bar bezahlt, bekommt 20 \% Rabatt. Kreuze an, wie viel Euro man dabei spart:
\pfkreuzzeile{100 €}\pfkreuzzeile{440 €}\pfkreuzzeile{110 €}\pfkreuzzeile{660 €}
\fusshilfe{35: 110 €}
\end{pfnr}
\begin{pfnr}{P10 ’15}{36.}{3 BE}
Für den Berliner Fernsehturm gelten diese Eintrittspreise: Erwachsene zahlen 12,00 €, Kinder von 3 bis 16 Jahren 7,50 €. In den Wintermonaten gibt es auf den Eintritt 20 \% Rabatt (Rabatt heißt Preisnachlass). Im Winter kommen Vater und Mutter mit ihrem 13-jährigen Sohn und dem Großvater zum Fernsehturm und nutzen den Winterrabatt. Entscheide für jeden Term, ob er den Rabatt richtig berechnet. Kreuze richtig oder falsch an:
\par\smallskip\noindent \begingroup\renewcommand{\arraystretch}{1.5}\begin{tabular}{|l|c|c|}\hline Term & richtig & falsch\\ \hline $(2\cdot 12\,€ + 2\cdot 7{,}50\,€)\cdot\tfrac{20}{100}$ & $\square$ & $\square$ \\ $\tfrac{1}{5}\cdot(36\,€ + 7{,}50\,€)$ & $\square$ & $\square$ \\ $\tfrac{20\cdot(12\,€ + 12\,€ + 12\,€ + 7{,}50\,€)}{100}$ & $\square$ & $\square$ \\ \hline\end{tabular}\endgroup\par
\pfzwfrage{Berechne den vollen Eintrittspreis der Familie ohne Rabatt.}
\fusshilfe{36: Term 1 falsch; Term 2 richtig; Term 3 richtig}
\end{pfnr}
\pfstufekopf[sprozentsatz]{Prozentsatz $p$}{in 1 der letzten 5 Prüfungen}
\pfgruppe{}
\begin{pfnr}{eigene Aufgabe}{37.}{}
Teile den Streifen in gleiche Teile ein. Jeder Teil soll $20\,\%$ sein.
\par\smallskip\noindent \streifenleer[0]\par
\fusshilfe{37: $5$ gleiche Teile}
\end{pfnr}
\begin{pfnr}{eigene Aufgabe}{38.}{}
Teile den Streifen in gleiche Teile ein. Jeder Teil soll $5\,\%$ sein.
\par\smallskip\noindent \streifenleer[0]\par
\fusshilfe{38: $20$ gleiche Teile}
\end{pfnr}
\begin{pfnr}{eigene Aufgabe}{39.}{}
$45$ Gäste essen vegetarisch. Eingeladen sind $150$ Gäste. Welche Zahl ist hier das Ganze?
\rechenplatz{3}
\fusshilfe{39: $150$ Gäste}
\end{pfnr}
\begin{pfnr}{eigene Aufgabe}{40.}{}
Ein Pullover kostete $60\,€$. Jetzt kostet er $45\,€$. Welcher Preis ist hier das Ganze?
\fusshilfe{40: der alte Preis, $60\,€$}
\end{pfnr}
\begin{pfnr}{eigene Aufgabe}{41.}{}
Der ganze Streifen steht für $50$ Kinder. Der graue Teil steht für $10$ Kinder. Wie viel Prozent des ganzen Streifens sind grau?
\par\smallskip\noindent \streifenfeld{20}{$0$}{$50$ Kinder}\par
\par\noindent \leerfeld[\%] \par
\fusshilfe{41: $20\,\%$}
\end{pfnr}
\begin{pfnr}{eigene Aufgabe}{42.}{}
Der ganze Streifen steht für $40$ kg. Der graue Teil steht für $30$ kg. Wie viel Prozent des ganzen Streifens sind grau?
\par\smallskip\noindent \streifenfeld[4]{75}{$0$}{$40$ kg}\par
\par\noindent \leerfeld[\%] \par
\fusshilfe{42: $75\,\%$}
\end{pfnr}
\pfgruppe{}
\begin{pfnr}{eigene Aufgabe}{43.}{}
Im Chor singen $14$ Mädchen und $6$ Jungen. Wie viel Prozent der Kinder im Chor sind Jungen?
\par\noindent \leerfeld[\%] \par
\rechenplatz{4}
\fusshilfe{43: $30\,\%$}
\end{pfnr}
\begin{pfnr}{eigene Aufgabe}{44.}{}
Eine Jacke kostete vorher $80\,€$. Jetzt kostet sie $60\,€$. Wie viel Prozent Rabatt gibt es auf den alten Preis?
\par\noindent \leerfeld[\%] \par
\rechenplatz{4}
\fusshilfe{44: $25\,\%$}
\end{pfnr}
\pfgruppe{}
\begin{pfnr}{eigene Aufgabe}{45.}{}
In einer Schule wurden $100$ Kinder befragt. $17$ davon spielen ein Instrument. Wie viel Prozent der befragten Kinder spielen ein Instrument?
\par\noindent \leerfeld[\%] \par
\rechenplatz{3}
\fusshilfe{45: $17\,\%$}
\end{pfnr}
\begin{pfnr}{eigene Aufgabe}{46.}{}
In einer Schule wurden $100$ Kinder befragt. $68$ davon spielen ein Instrument. Wie viel Prozent der befragten Kinder spielen ein Instrument?
\par\noindent \leerfeld[\%] \par
\rechenplatz{3}
\fusshilfe{46: $68\,\%$}
\end{pfnr}
\begin{pfnr}{eigene Aufgabe}{47.}{}
Eine Bäckerei wirft abends Brötchen weg. Am Montag sind es $36$ von $240$ Brötchen. Am Dienstag sind es $27$ von $150$ Brötchen. An welchem Tag war der Anteil weggeworfener Brötchen größer?
\rechenplatz{4}
\fusshilfe{47: am Dienstag}
\end{pfnr}
\pfab{$p = W : G$}
\begin{pfnr}{eigene Aufgabe}{48.}{}
$17$ von $24$ Kindern waren im Schwimmbad. Wie viel Prozent der Kinder waren im Schwimmbad? Runde auf eine Stelle nach dem Komma.
\par\noindent \leerfeld[\%] \par
\rechenplatz{3}
\fusshilfe{48: $\tfrac{17}{24} \approx 70{,}8\,\%$}
\end{pfnr}
\begin{pfnr}{P10 ’18}{49.}{1 BE}
Für ein Fest backt Eva 16 Pfannkuchen: 14 füllt sie mit Marmelade, 2 mit Senf. Gib an, wie viel Prozent der Pfannkuchen eine Senffüllung haben.
\rechenplatz{2}
\fusshilfe{49: 12,5 \%}
\end{pfnr}
\begin{pfnr}{P10 ’15}{50.}{4 BE}
In einer Saison der 2. Fußball-Bundesliga wurden 83 Elfmeter gepfiffen; 71 davon führten zu einem Tor. Berechne, wie viel Prozent der Elfmeter nicht verwandelt wurden. Zeichne den Anteil als Sektor in den vorbereiteten Kreis ein. Gib an, wie groß der Mittelpunktswinkel dieses Sektors ist.
\par\smallskip\noindent \kreisleer[2]\par
\pfzwfrage{Berechne, wie viele Elfmeter nicht verwandelt wurden.}
\pfzwfrage{Berechne den Anteil dieser Elfmeter an allen Elfmetern.}
\pfzwfrage{Berechne den Mittelpunktswinkel, der zu diesem Anteil gehört.}
\rechenplatz{4}
\fusshilfe{50: $\tfrac{12}{83}$ $\approx$ 14,5 \%; Sektor $\approx$ 52°; $\approx$ 52°}
\end{pfnr}
\pfstufekopf[sgrundwert]{Grundwert $G$}{in 2 der letzten 5 Prüfungen}
\pfgruppe{}
\begin{pfnr}{eigene Aufgabe}{51.}{}
Der graue Abschnitt ist $10\,\%$ des Ganzen. Trage den Abschnitt so oft ab, bis du bei $100\,\%$ bist. Markiere das Ganze.
\par\smallskip\noindent \streifen[0]{10}{$0\,\%$}{}\par
\par\noindent \leerfeld[Abschnitte] \par
\fusshilfe{51: $10$ Abschnitte bis $100\,\%$}
\end{pfnr}
\begin{pfnr}{eigene Aufgabe}{52.}{}
Der graue Abschnitt ist $50\,\%$ des Ganzen. Trage den Abschnitt so oft ab, bis du bei $100\,\%$ bist. Markiere das Ganze.
\par\smallskip\noindent \streifen[0]{50}{$0\,\%$}{}\par
\par\noindent \leerfeld[Abschnitte] \par
\fusshilfe{52: $2$ Abschnitte bis $100\,\%$}
\end{pfnr}
\begin{pfnr}{eigene Aufgabe}{53.}{}
In einer Aufgabe steht: „Von $240$ Befragten sagen $60$ ja. Wie viel Prozent sind das?“ Kreuze an, was gesucht ist.
\pfkreuzzeile{Prozentwert}\pfkreuzzeile{Prozentsatz}\pfkreuzzeile{Grundwert}
\fusshilfe{53: Prozentsatz}
\end{pfnr}
\begin{pfnr}{eigene Aufgabe}{54.}{}
In einer Aufgabe steht: „$30\,\%$ der Klasse sind $9$ Kinder. Wie viele Kinder hat die Klasse?“ Kreuze an, was gesucht ist.
\pfkreuzzeile{Prozentwert}\pfkreuzzeile{Prozentsatz}\pfkreuzzeile{Grundwert}
\fusshilfe{54: Grundwert}
\end{pfnr}
\pfgruppe{}
\begin{pfnr}{eigene Aufgabe}{55.}{}
Ein Händler hat $24$ kg Äpfel verkauft. Das sind $10\,\%$ seiner Äpfel. Wie viel kg Äpfel hatte er?
\par\noindent \leerfeld[kg] \par
\rechenplatz{4}
\fusshilfe{55: $240$ kg}
\end{pfnr}
\begin{pfnr}{eigene Aufgabe}{56.}{}
Ein Händler hat $24$ kg Äpfel verkauft. Das sind $50\,\%$ seiner Äpfel. Wie viel kg Äpfel hatte er?
\par\noindent \leerfeld[kg] \par
\rechenplatz{4}
\fusshilfe{56: $48$ kg}
\end{pfnr}
\begin{pfnr}{eigene Aufgabe}{57.}{}
Im Tank sind noch $12$ l. Das sind $20\,\%$ der Tankfüllung. Wie viele Liter passen in den Tank?
\par\noindent \leerfeld[l] \par
\rechenplatz{3}
\fusshilfe{57: $60$ l}
\end{pfnr}
\pfab{$G = W : p$}
\begin{pfnr}{eigene Aufgabe}{58.}{}
$36\,\%$ einer Menge sind $81$ kg. Wie viel kg ist die ganze Menge?
\par\noindent \leerfeld[kg] \par
\rechenplatz{4}
\fusshilfe{58: $225$ kg}
\end{pfnr}
\pfgruppe{}
\begin{pfnr}{eigene Aufgabe}{59.}{}
Tims Handy zeigt $28\,\%$ Akku. Laut Anzeige reicht das noch für $7$ Stunden Musik. Wie lange hält ein voller Akku? Reicht ein voller Akku für eine Busfahrt von $22$ Stunden?
\rechenplatz{4}
\fusshilfe{59: Ja}
\end{pfnr}
\pfgruppe{}
\begin{pfnr}{P10 ’25}{60.}{1 BE}
Ein T-Shirt ist um 6 € billiger geworden. Dieser Nachlass entspricht 20 \% des ursprünglichen Preises. Gib den ursprünglichen Preis an.
\rechenplatz{2}
\fusshilfe{60: 30 €, Tipp: $G$ = 6 € : 0,2}
\end{pfnr}
\begin{pfnr}{P10 ’23}{61.}{1 BE}
Mia gewinnt bei einem Wettbewerb Geld. Sie gibt 25 \% des Gewinns an ihre Eltern ab; das sind 200 €. Gib an, wie hoch der gesamte Gewinn war.
\rechenplatz{2}
\fusshilfe{61: 800 €, Tipp: $G$ = 200 € : 0,25}
\end{pfnr}
\pfstufekopf[serhhungundvernderung]{Erhöhung und Veränderung in Prozent}{in 3 der letzten 5 Prüfungen}
\pfgruppe{}
\begin{pfnr}{eigene Aufgabe}{62.}{}
Ein Preis wird um $20\,\%$ gesenkt. Kreuze an, was dasselbe bedeutet.
\pfkreuzzeile{auf $20\,\%$ gesenkt}\pfkreuzzeile{auf $80\,\%$ gesenkt}\pfkreuzzeile{um $80\,\%$ gesenkt}
\fusshilfe{62: auf $80\,\%$ gesenkt}
\end{pfnr}
\begin{pfnr}{eigene Aufgabe}{63.}{}
Jetzt hat der Verein $150$ Mitglieder. Im letzten Jahr waren es $120$. Welche Zahl ist der alte Wert?
\rechenplatz{3}
\fusshilfe{63: $120$ Mitglieder}
\end{pfnr}
\begin{pfnr}{eigene Aufgabe}{64.}{}
Die Miete wurde auf $690\,€$ erhöht. Vorher zahlte Familie Berg $600\,€$. Welcher Betrag ist der alte Wert?
\fusshilfe{64: $600\,€$}
\end{pfnr}
\begin{pfnr}{eigene Aufgabe}{65.}{}
Ein Preis wird auf $60\,\%$ gesenkt. Kreuze an, was dasselbe bedeutet.
\pfkreuzzeile{um $60\,\%$ gesenkt}\pfkreuzzeile{um $40\,\%$ gesenkt}\pfkreuzzeile{auf $40\,\%$ gesenkt}
\fusshilfe{65: um $40\,\%$ gesenkt}
\end{pfnr}
\pfgruppe{}
\begin{pfnr}{eigene Aufgabe}{66.}{}
Ein Preis von $70\,€$ wird um $10\,\%$ erhöht. Der Streifen zeigt die $10\,\%$. Wie hoch ist der neue Preis?
\par\smallskip\noindent \streifen[10]{10}{$0$}{$70\,€$}\par
\par\noindent \leerfeld[€] \par
\rechenplatz{4}
\fusshilfe{66: $77\,€$}
\end{pfnr}
\begin{pfnr}{eigene Aufgabe}{67.}{}
Die Monatskarte wird um $10\,\%$ teurer. Jetzt kostet sie $66\,€$. Wie hoch war der alte Preis?
\par\noindent \leerfeld[€] \par
\rechenplatz{3}
\fusshilfe{67: $60\,€$}
\end{pfnr}
\begin{pfnr}{eigene Aufgabe}{68.}{}
Ein Preis von $70\,€$ wird um $50\,\%$ erhöht. Der Streifen zeigt die $50\,\%$. Wie hoch ist der neue Preis?
\par\smallskip\noindent \streifen[2]{50}{$0$}{$70\,€$}\par
\par\noindent \leerfeld[€] \par
\rechenplatz{4}
\fusshilfe{68: $105\,€$}
\end{pfnr}
\begin{pfnr}{eigene Aufgabe}{69.}{}
In einer AG steigt die Zahl der Kinder von $40$ auf $50$. Um wie viel Prozent ist die Zahl gestiegen?
\par\noindent \leerfeld[\%] \par
\rechenplatz{3}
\fusshilfe{69: $25\,\%$}
\end{pfnr}
\begin{pfnr}{eigene Aufgabe}{70.}{}
Ein Buch kostet mit $7\,\%$ Mehrwertsteuer (brutto) $53{,}50\,€$. Wie viel kostet es ohne Mehrwertsteuer (netto)?
\par\noindent \leerfeld[€] \par
\rechenplatz{3}
\fusshilfe{70: $50\,€$}
\end{pfnr}
\begin{pfnr}{eigene Aufgabe}{71.}{2 BE}
Ein Baum war $4{,}20$ m hoch, jetzt ist er $4{,}83$ m hoch. Um wie viel Prozent ist er gewachsen?
\par\noindent \leerfeld[\%] \par
\rechenplatz{3}
\fusshilfe{71: $15\,\%$}
\end{pfnr}
\pfgruppe{}
\begin{pfnr}{eigene Aufgabe}{72.}{}
Ein Fahrrad kostet $600\,€$. Bei Angebot A gibt es $10\,\%$ Rabatt. Auf den neuen Preis gibt es noch einmal $10\,\%$. Lisa hat $490\,€$ gespart. Reicht ihr Geld für Angebot A?
\rechenplatz{6}
\fusshilfe{72: Ja}
\end{pfnr}
\begin{pfnr}{P10 ’20}{73.}{2 BE}
Das Säulendiagramm zeigt die Besucherzahlen eines Freibads an den Tagen der ersten Ferienwoche. Die beiden folgenden Aussagen sind richtig: „Die Besucherzahlen der ersten Ferienwoche haben eine Spannweite von 150.“ „Am Sonntag gab es ein Drittel mehr Besucher als am Mittwoch.“ Weise beide Aussagen rechnerisch nach.
\par\smallskip\noindent \begin{tikzpicture}\begin{axis}[ybar,width=11.1cm,height=4.6cm,ymin=0,ymax=336.00,symbolic x coords={Mo,Di,Mi,Do,Fr,Sa,So},xtick=data,enlarge x limits=0.08,bar width=6mm,ylabel={Besucher je Wochentag},ylabel style={font=\small},tick label style={font=\small},nodes near coords,every node near coord/.append style={font=\scriptsize,/pgf/number format/.cd,use comma,1000 sep={}},ymajorgrids,grid style={mbgitter}]\addplot[fill=mbkasten,draw=black] coordinates {(Mo,170.0) (Di,130.0) (Mi,210.0) (Do,180.0) (Fr,190.0) (Sa,240.0) (So,280.0)};\end{axis}\end{tikzpicture}\par
\pfzwfrage{Berechne ein Drittel der Besucherzahl vom Mittwoch.}
\rechenplatz{4}
\fusshilfe{73: 280 $-$ 130 = 150; 210 + 210 : 3 = 280}
\end{pfnr}
\begin{pfnr}{P10 ’20}{\llap{\textasteriskcentered\,}74.}{3 BE}
Nach einer Studie lag der Wasserverbrauch 2019 bei durchschnittlich 54 Tausend Litern pro Person und Jahr; er soll jedes Jahr um 3 \% steigen. Diese Vorhersage beschreibt die Gleichung $y = 54\cdot 1{,}03^{x}$ ($x$: Jahre ab 2019, $y$: Verbrauch pro Person in Tausend Litern). Eine zweite Studie erwartet, dass der Verbrauch pro Person von 2019 bis 2033 insgesamt um etwa 40 \% zunimmt. Untersuche, ob beide Studien für 2033 denselben Verbrauch vorhersagen.
\pfzwfrage{Berechne mit der Gleichung den Verbrauch pro Person im Jahr 2033.}
\rechenplatz{4}
\end{pfnr}
\pfgruppe{}
\begin{pfnr}{eigene Aufgabe}{75.}{}
Ein Zinssatz steigt von $2\,\%$ auf $3\,\%$. Um wie viele Prozentpunkte ist er gestiegen? Um wie viel Prozent ist er gestiegen?
\rechenplatz{4}
\fusshilfe{75: $50\,\%$ mehr}
\end{pfnr}
\begin{pfnr}{P10 ’26}{76.}{2 BE}
An einer Tankstelle wurde eine Woche lang täglich um 18 Uhr der Benzinpreis notiert (€ pro Liter): Mo 1,77, Di 1,78, Mi 1,77, Do 1,79, Fr 1,84, Sa 1,82, So 1,85. Berechne die prozentuale Preissteigerung von Donnerstag auf Freitag.
\rechenplatz{3}
\fusshilfe{76: $\tfrac{5}{179}$ $\approx$ 2,8 \%}
\end{pfnr}
\begin{pfnr}{P10 ’24}{77.}{1 BE}
Ein Schwimmbad verlangt bisher 3,50 € Eintritt. Der Preis steigt um 20 \%. Gib den neuen Eintrittspreis in Euro an.
\rechenplatz{2}
\fusshilfe{77: 4,20 €}
\end{pfnr}
\begin{pfnr}{P10 ’16}{78.}{2 BE}
Das Säulendiagramm zeigt, wie viele Besucher ein Filmpark in den Jahren 2007 bis 2015 hatte. Berechne, um wie viel Prozent die Besucherzahl von 2009 auf 2010 zugenommen hat.
\par\smallskip\noindent \begin{tikzpicture}\begin{axis}[ybar,width=13.7cm,height=4.6cm,ymin=0,ymax=720.00,symbolic x coords={2007,2008,2009,2010,2011,2012,2013,2014,2015},xtick=data,enlarge x limits=0.08,bar width=6mm,ylabel={Besucher in Tausend},ylabel style={font=\small},tick label style={font=\small},nodes near coords,every node near coord/.append style={font=\scriptsize,/pgf/number format/.cd,use comma,1000 sep={}},ymajorgrids,grid style={mbgitter}]\addplot[fill=mbkasten,draw=black] coordinates {(2007,600.0) (2008,420.0) (2009,320.0) (2010,400.0) (2011,260.0) (2012,260.0) (2013,275.0) (2014,300.0) (2015,250.0)};\end{axis}\end{tikzpicture}\par
\rechenplatz{3}
\fusshilfe{78: 25 \%}
\end{pfnr}
\begin{pfnr}{P10 ’22}{79.}{2 BE}
Das Säulendiagramm zeigt, wie viele von jeweils 1200 befragten Jugendlichen in den Jahren 2017 bis 2021 mehrmals pro Woche ein Buch lesen. Ermittle, um wie viel Prozent diese Zahl von 2020 auf 2021 zurückgegangen ist.
\par\smallskip\noindent \begin{tikzpicture}\begin{axis}[ybar,width=8.5cm,height=4.6cm,ymin=390,ymax=508.75,symbolic x coords={2017,2018,2019,2020,2021},xtick=data,enlarge x limits=0.08,bar width=6mm,ylabel={Anzahl},ylabel style={font=\small},tick label style={font=\small},nodes near coords,every node near coord/.append style={font=\scriptsize,/pgf/number format/.cd,use comma,1000 sep={}},ymajorgrids,grid style={mbgitter}]\addplot[fill=mbkasten,draw=black] coordinates {(2017,468.0) (2018,432.0) (2019,485.0) (2020,470.0) (2021,400.0)};\end{axis}\end{tikzpicture}\par
\rechenplatz{3}
\fusshilfe{79: $\tfrac{7}{47}$ $\approx$ 14,9 \%}
\end{pfnr}
\begin{pfnr}{P10 ’15}{80.}{2 BE}
Für den Berliner Fernsehturm gelten diese Eintrittspreise: Erwachsene zahlen 12,00 €, Kinder von 3 bis 16 Jahren 7,50 €. In den Wintermonaten gibt es auf den Eintritt 20 \% Rabatt (Rabatt heißt Preisnachlass). Berechne den Betrag in Euro, den diese Familie insgesamt für den Eintritt bezahlt.
\rechenplatz{3}
\fusshilfe{80: 34,80 €}
\end{pfnr}
\begin{pfnr}{P10 ’21}{81.}{5 BE}
Das Säulendiagramm gibt an, wie viele Container (in Millionen) in den Jahren 2009 bis 2017 im Hamburger Hafen be- und entladen wurden. Berechne die Spannweite und den Durchschnitt (das arithmetische Mittel) dieser Containerzahlen. Berechne, um wie viel Prozent die Containerzahl von 2010 auf 2011 gestiegen ist.
\par\smallskip\noindent \begin{tikzpicture}\begin{axis}[ybar,width=13.7cm,height=4.6cm,ymin=0,ymax=11.64,symbolic x coords={2009,2010,2011,2012,2013,2014,2015,2016,2017},xtick=data,enlarge x limits=0.08,bar width=6mm,ylabel={Container in Millionen},ylabel style={font=\small},tick label style={font=\small},nodes near coords,every node near coord/.append style={font=\scriptsize,/pgf/number format/.cd,use comma,1000 sep={}},ymajorgrids,grid style={mbgitter}]\addplot[fill=mbkasten,draw=black] coordinates {(2009,7.0) (2010,7.9) (2011,9.0) (2012,8.9) (2013,9.3) (2014,9.7) (2015,8.8) (2016,8.9) (2017,8.8)};\end{axis}\end{tikzpicture}\par
\rechenplatz{5}
\fusshilfe{81: 2,7 Mio.; 8,7 Mio.; $\tfrac{11}{79}$ $\approx$ 13,9 \%}
\end{pfnr}
\pfstufekopf[saussagenprfen]{Aussagen prüfen}{in 2 der letzten 5 Prüfungen}
\pfgruppe{}
\begin{pfnr}{eigene Aufgabe}{82.}{}
Der Preis stieg von $2{,}40\,€$ auf $2{,}70\,€$. Welcher Preis ist der alte Preis?
\rechenplatz{2}
\fusshilfe{82: $2{,}40\,€$}
\end{pfnr}
\begin{pfnr}{eigene Aufgabe}{83.}{}
Nach dem Rabatt kostet das Rad $340\,€$. Vorher kostete es $425\,€$. Welcher Preis ist der alte Preis?
\rechenplatz{3}
\fusshilfe{83: $425\,€$}
\end{pfnr}
\begin{pfnr}{eigene Aufgabe}{84.}{}
Ein Preis wird um $15\,\%$ erhöht. Kreuze an, was dasselbe bedeutet.
\pfkreuzzeile{auf $115\,\%$ erhöht}\pfkreuzzeile{auf $15\,\%$ erhöht}\pfkreuzzeile{auf $85\,\%$ erhöht}
\fusshilfe{84: auf $115\,\%$ erhöht}
\end{pfnr}
\begin{pfnr}{eigene Aufgabe}{85.}{}
Ein Preis wird auf $125\,\%$ erhöht. Kreuze an, was dasselbe bedeutet.
\pfkreuzzeile{um $125\,\%$ erhöht}\pfkreuzzeile{um $75\,\%$ erhöht}\pfkreuzzeile{um $25\,\%$ erhöht}
\fusshilfe{85: um $25\,\%$ erhöht}
\end{pfnr}
\pfgruppe{}
\begin{pfnr}{eigene Aufgabe}{86.}{}
Ein Preis von $70\,€$ wird um $25\,\%$ gesenkt. Der Streifen zeigt die $25\,\%$. Wie hoch ist der neue Preis?
\par\smallskip\noindent \streifen[4]{25}{$0$}{$70\,€$}\par
\par\noindent \leerfeld[€] \par
\rechenplatz{4}
\fusshilfe{86: $52{,}50\,€$}
\end{pfnr}
\begin{pfnr}{eigene Aufgabe}{87.}{}
Ein Preis von $70\,€$ wird um $20\,\%$ gesenkt. Der Streifen zeigt die $20\,\%$. Wie hoch ist der neue Preis?
\par\smallskip\noindent \streifen[5]{20}{$0$}{$70\,€$}\par
\par\noindent \leerfeld[€] \par
\rechenplatz{4}
\fusshilfe{87: $56\,€$}
\end{pfnr}
\begin{pfnr}{eigene Aufgabe}{88.}{}
Eine Straße steigt auf $200$ m waagerechter Strecke um $14$ m. Wie groß ist die Steigung in Prozent?
\par\noindent \leerfeld[\%] \par
\rechenplatz{3}
\fusshilfe{88: $7\,\%$}
\end{pfnr}
\begin{pfnr}{eigene Aufgabe}{89.}{}
Der Streifen zeigt ein Ganzes. Ein Teil davon ist grau. Wie viel Prozent des ganzen Streifens sind grau?
\par\smallskip\noindent \streifenfeld{70}{$0\,\%$}{$100\,\%$}\par
\par\noindent \leerfeld[\%] \par
\rechenplatz{2}
\fusshilfe{89: $70\,\%$}
\end{pfnr}
\begin{pfnr}{eigene Aufgabe}{90.}{}
Der Streifen zeigt ein Ganzes. Ein Teil davon ist grau. Wie viel Prozent des ganzen Streifens sind grau?
\par\smallskip\noindent \streifenfeld{30}{$0\,\%$}{$100\,\%$}\par
\par\noindent \leerfeld[\%] \par
\rechenplatz{2}
\fusshilfe{90: $30\,\%$}
\end{pfnr}
\pfgruppe{}
\begin{pfnr}{eigene Aufgabe}{91.}{}
Die Klassen 7a und 7b stimmen über den Wandertag ab. In der 7a sind $18$ von $25$ Kindern für den Zoo. In der 7b sind $21$ von $30$ Kindern für den Zoo. In welcher Klasse ist der Anteil für den Zoo größer?
\rechenplatz{4}
\fusshilfe{91: in der 7a ist der Anteil größer}
\end{pfnr}
\pfgruppe{}
\begin{pfnr}{P10 ’19}{92.}{2 BE}
Die Tabelle gibt an, welcher Anteil der Mädchen und der Jungen ein Medium in der Freizeit nutzt (Mädchen / Jungen): Smartphone 98 \% / 95 \%, Fernseher 81 \% / 75 \%, Konsolenspiele 14 \% / 72 \%, Bücher 46 \% / 30 \%. Zu der Tabelle gibt es zwei Aussagen: Kreuze die falsche Aussage an. Berichtige diese Aussage.
\pfkreuzzeile{„Drei Viertel der Mädchen sehen in ihrer Freizeit fern.“}\pfkreuzzeile{„Bei Jungen sind Konsolenspiele deutlich beliebter als bei Mädchen.“}
\rechenplatz{2}
\end{pfnr}
\pfgruppe{}
\begin{pfnr}{P10 ’25}{93.}{3 BE}
Eine Rampe für Rollstühle darf höchstens 6 \% Steigung haben, damit man sie ohne fremde Hilfe hinauffahren kann. Vervollständige den Satz: „Bei 6 \% Steigung steigt man auf \_\_\_ Metern waagerechter Strecke um \_\_\_ Meter an.“ Am Hauseingang soll eine Rampe eine 16 cm hohe Stufe überwinden; waagerecht ist sie 170 cm lang. Frau Yücel meint: „Unsere Rampe wird zu steil, ihre Steigung liegt über 6 \%.“ Überprüfe, ob Frau Yücel recht hat.
\par\smallskip\noindent \begin{tikzpicture}[line width=0.6pt]\draw (0,0) -- (6,0) -- (6,1.4) -- cycle;\draw (5.7,0) -- (5.7,0.3) -- (6,0.3);\node[below,font=\small] at (3,0) {170 cm};\node[right,font=\small] at (6,0.7) {16 cm};\node[above left,font=\small] at (3,0.7) {$x$};\node[font=\small] at (1.1,0.12) {$\alpha$};\node[font=\scriptsize,mbgrau,anchor=west] at (0,-0.75) {nicht maßstabsgerecht};\end{tikzpicture}\par
\pfzwfrage{Berechne das Verhältnis von Höhe zu waagerechter Länge der geplanten Rampe.}
\rechenplatz{3}
\end{pfnr}
\begin{pfnr}{P10 ’17}{\llap{\textasteriskcentered\,}94.}{2 BE}
Von allen Kindern, die 2013 in Deutschland zur Welt kamen, waren 49,4 \% das erste Kind ihrer Mutter, 34,4 \% das zweite, 11,2 \% das dritte und 5,0 \% das vierte oder ein späteres. Eine Forscherin folgert daraus: „Über die Hälfte dieser Kinder hat mindestens ein Geschwisterkind.“ Gib eine Begründung für diese Folgerung an.
\rechenplatz{2}
\fusshilfe{94: 34,4 \% + 11,2 \% + 5,0 \% = 50,6 \% > 50 \%}
\end{pfnr}
\pfstufekopf[sprozentauseinerberec]{Prozent aus einer berechneten Fläche}{in 1 der letzten 5 Prüfungen}
\pfgruppe{}
\begin{pfnr}{eigene Aufgabe}{\llap{\textasteriskcentered\,}95.}{}
Aus einer quadratischen Korkplatte mit $40$ cm Seitenlänge werden kreisrunde Untersetzer mit dem Radius $5$ cm ausgeschnitten, so viele wie möglich, in Reihen nebeneinander. Berechne, wie viel Prozent der Platte Abfall sind. Runde auf eine Stelle nach dem Komma.
\rechenplatz{6}
\fusshilfe{95: $\approx 21{,}5\,\%$}
\end{pfnr}
\pfgruppe{}
\begin{pfnr}{P10 ’23}{\llap{\textasteriskcentered\,}96.}{4 BE}
Für Konservendosen werden kreisrunde Deckel mit dem Radius 4 cm aus Blechstreifen ausgestanzt, die 8 cm breit und 1 m lang sind. Aus jedem Streifen sollen möglichst viele Deckel entstehen. Der Abfall darf höchstens 25 \% eines Streifens ausmachen. Überprüfe rechnerisch, ob diese Vorgabe eingehalten wird.
\pfzwfrage{Berechne, wie viele Deckel aus einem Streifen ausgestanzt werden können.}
\pfzwfrage{Berechne die Fläche aller Deckel eines Streifens.}
\pfzwfrage{Berechne die Fläche eines Streifens.}
\rechenplatz{5}
\end{pfnr}
\pfstufekopf[szinsenundzinssatz]{Zinsen $Z$ und Zinssatz $p$}{selten geprüft}
\pfgruppe{}
\begin{pfnr}{eigene Aufgabe}{97.}{}
Auf einem Konto liegen 2\,400\,€. Nach einem Jahr gab es 48\,€ Zinsen. Die Frage lautet: „Wie hoch ist der Zinssatz?“ Kreuze an, was gesucht ist. Rechne nicht. 
\pfkreuzzeile{das Kapital (Grundwert)}\pfkreuzzeile{der Zinssatz (Prozentsatz)}\pfkreuzzeile{die Zinsen (Prozentwert)}
\fusshilfe{97: der Zinssatz}
\end{pfnr}
\begin{pfnr}{eigene Aufgabe}{98.}{}
In der Prozentrechnung gilt $W = G \cdot p\,\%$. Welche Formel gehört dazu in der Zinsrechnung? Kreuze an. Rechne nicht. 
\pfkreuzzeile{$K = Z \cdot p\,\%$}\pfkreuzzeile{$Z = K \cdot p\,\%$}\pfkreuzzeile{$p\,\% = K \cdot Z$}
\fusshilfe{98: $K \cdot p\,\%$}
\end{pfnr}
\pfgruppe{}
\begin{pfnr}{eigene Aufgabe}{99.}{}
900\,€ Kapital bringen in einem Jahr 27\,€ Zinsen. Wie hoch ist der Zinssatz?
\par\noindent \leerfeld[\%] \par
\rechenplatz{3}
\fusshilfe{99: 3\,\%}
\end{pfnr}
\begin{pfnr}{eigene Aufgabe}{100.}{}
Auf einem Konto liegen 4\,000\,€. In einem Jahr gibt es dafür 60\,€ Zinsen. Wie hoch ist der Zinssatz?
\par\noindent \leerfeld[\%] \par
\rechenplatz{3}
\fusshilfe{100: 1,5\,\%}
\end{pfnr}
\pfab{$Z = K \cdot p$}
\begin{pfnr}{eigene Aufgabe}{101.}{}
Die Zinsen für ein Jahr sind 21\,€. Das sind 3\,\% vom Kapital. Wie viel Euro ist das Kapital?
\par\noindent \leerfeld[€] \par
\rechenplatz{3}
\fusshilfe{101: 700\,€}
\end{pfnr}
\begin{pfnr}{eigene Aufgabe}{102.}{}
Weise mit der Tabelle nach, dass die Bank 3\,\% Zinsen im Jahr zahlt.
\par\smallskip\noindent \sachtabelle{lrrr}{Buchung & Einzahlung & Zinsen & Guthaben}{Jahresbeginn & 300,00\,€ & – & 300,00\,€\\ Jahresende & – & 9,00\,€ & 309,00\,€}\par
\rechenplatz{4}
\fusshilfe{102: stimmt}
\end{pfnr}
\pfgruppe{}
\begin{pfnr}{eigene Aufgabe}{103.}{}
Auf einem Sparbuch liegen 300\,€. Die Bank zahlt 2\,\% Zinsen im Jahr. Wie viel Euro Zinsen gibt es nach einem Jahr?
\par\noindent \leerfeld[€] \par
\rechenplatz{3}
\fusshilfe{103: 6\,€}
\end{pfnr}
\begin{pfnr}{eigene Aufgabe}{104.}{}
Auf einem Sparbuch liegen 700\,€. Die Bank zahlt 3\,\% Zinsen im Jahr. Wie viel Euro Zinsen gibt es nach einem Jahr?
\par\noindent \leerfeld[€] \par
\rechenplatz{3}
\fusshilfe{104: 21\,€}
\end{pfnr}
\begin{pfnr}{P10 ’15}{105.}{1 BE}
Auf einem Konto liegen 400 €, die mit 2 \% im Jahr verzinst werden. Gib an, wie viel Euro Zinsen nach einem Jahr dazukommen.
\rechenplatz{2}
\fusshilfe{105: 8 €}
\end{pfnr}
\begin{pfnr}{P10 ’14}{106.}{1 BE}
Auf einem Sparkonto liegen 10 000 €. Nach einem Jahr werden 230 € Zinsen gutgeschrieben. Bestimme den Zinssatz.
\rechenplatz{2}
\fusshilfe{106: 2,3 \%}
\end{pfnr}
\pfgruppe{}
\begin{pfnr}{P10 ’14}{107.}{1 BE}
Lena zahlt seit ihrem 11. Geburtstag jedes Jahr an ihrem Geburtstag 200,00 € auf ein Sparbuch ein. Am 15. Geburtstag hält sie in einer Tabelle fest, wie sich ihr Guthaben entwickelt hat. Das Guthaben wird mit 2 \% pro Jahr verzinst. Weise dies anhand der Werte für das erste Jahr nach.
\par\smallskip\noindent \sachtabelle{lrrrrr}{Geburtstag & 11. & 12. & 13. & 14. & 15.}{Zinsen & – & 4,00 € & 8,08 € & 12,24 € & \leerzelle\\ Einzahlung & 200,00 € & 200,00 € & 200,00 € & 200,00 € & 200,00 €\\ Guthaben & 200,00 € & 404,00 € & 612,08 € & \leerzelle & 1040,81 €}\par
\rechenplatz{2}
\fusshilfe{107: 200 € · 0,02 = 4,00 €}
\end{pfnr}
\pfstufekopf[szinseszinsundguthabe]{Zinseszins und Guthabentabelle}{selten geprüft}
\pfgruppe{}
\begin{pfnr}{eigene Aufgabe}{108.}{}
Auf einem Konto liegen am Anfang 600\,€. Nach einem Jahr sind es mit Zinsen 612\,€. Kreuze an, von welchem Betrag die Zinsen im zweiten Jahr berechnet werden. Rechne nicht. 
\pfkreuzzeile{vom neuen Guthaben 612\,€}\pfkreuzzeile{vom Startguthaben 600\,€}
\fusshilfe{108: vom neuen Guthaben 612\,€}
\end{pfnr}
\begin{pfnr}{eigene Aufgabe}{109.}{}
Ein Sparbuch startet mit 2\,000\,€. Nach einem Jahr stehen 2\,060\,€ darauf. Kreuze an, von welchem Betrag die Zinsen im zweiten Jahr berechnet werden. Rechne nicht. 
\pfkreuzzeile{vom Startguthaben 2\,000\,€}\pfkreuzzeile{vom neuen Guthaben 2\,060\,€}
\fusshilfe{109: vom neuen Guthaben 2\,060\,€}
\end{pfnr}
\pfgruppe{}
\begin{pfnr}{eigene Aufgabe}{110.}{}
Auf einem Konto liegen 5\,000\,€. Die Bank zahlt 1\,\% Zinsen im Jahr. Am Jahresende kommen die Zinsen auf das Konto. Wie hoch ist dann das neue Guthaben?
\par\noindent \leerfeld[€] \par
\rechenplatz{3}
\fusshilfe{110: 5\,050\,€}
\end{pfnr}
\begin{pfnr}{eigene Aufgabe}{111.}{}
Auf einem Konto liegen 2\,200\,€. Die Bank zahlt 5\,\% Zinsen im Jahr. Am Jahresende kommen die Zinsen auf das Konto. Wie hoch ist dann das neue Guthaben?
\par\noindent \leerfeld[€] \par
\rechenplatz{3}
\fusshilfe{111: 2\,310\,€}
\end{pfnr}
\begin{pfnr}{eigene Aufgabe}{112.}{}
Die Tabelle zeigt ein Sparkonto mit 2\,\% Zinsen im Jahr. Jedes Jahr werden 120\,€ eingezahlt. Ergänze die zwei leeren Felder. Kontrolliere dein Ergebnis mit dem Guthaben in Jahr 3.
\par\smallskip\noindent \sachtabelle{lrrr}{Jahr & Zinsen & Einzahlung & Guthaben}{1 & – & – & 350,00\,€\\ 2 & 7,00\,€ & 120,00\,€ & \leerzelle\\ 3 & \leerzelle & 120,00\,€ & 606,54\,€}\par
\par\noindent Guthaben \leerfeld[€] , Zinsen \leerfeld[€] \par
\rechenplatz{2}
\fusshilfe{112: Guthaben = 477,00\,€, Zinsen = 9,54\,€}
\end{pfnr}
\begin{pfnr}{eigene Aufgabe}{113.}{}
Auf einem Konto liegen am Start 600\,€. Der Zinssatz ist 2\,\%. Am Ende jedes Jahres kommen erst die Zinsen dazu, dann 120\,€ Einzahlung. Fülle die Tabelle aus. Wie hoch ist das Guthaben nach 3 Jahren?
\par\smallskip\noindent \sachtabelle{lrrr}{Jahr & Zinsen & Einzahlung & Guthaben}{Start & – & – & 600,00\,€\\ 1 & \leerzelle & 120,00\,€ & \leerzelle\\ 2 & \leerzelle & 120,00\,€ & \leerzelle\\ 3 & \leerzelle & 120,00\,€ & \leerzelle}\par
\par\noindent \leerfeld[€] \par
\rechenplatz{5}
\fusshilfe{113: 1\,003,97\,€}
\end{pfnr}
\begin{pfnr}{eigene Aufgabe}{114.}{}
Auf einem Konto liegen 3\,500\,€ für 4 Jahre. Die Bank zahlt 3\,\% Zinsen mit Zinseszins. Wie hoch ist das Guthaben am Ende? Runde auf Cent.
\par\noindent \leerfeld[€] \par
\rechenplatz{3}
\fusshilfe{114: $\approx$ 3\,939,28\,€}
\end{pfnr}
\begin{pfnr}{P10 ’14}{115.}{2 BE}
Lena zahlt seit ihrem 11. Geburtstag jedes Jahr an ihrem Geburtstag 200,00 € auf ein Sparbuch ein. Am 15. Geburtstag hält sie in einer Tabelle fest, wie sich ihr Guthaben entwickelt hat. Vervollständige die fehlenden Einträge der Tabelle.
\par\smallskip\noindent \sachtabelle{lrrrrr}{Geburtstag & 11. & 12. & 13. & 14. & 15.}{Zinsen & – & 4,00 € & 8,08 € & 12,24 € & \leerzelle\\ Einzahlung & 200,00 € & 200,00 € & 200,00 € & 200,00 € & 200,00 €\\ Guthaben & 200,00 € & 404,00 € & 612,08 € & \leerzelle & 1040,81 €}\par
\rechenplatz{3}
\fusshilfe{115: 824,32 €; 16,49 €}
\end{pfnr}
\begin{pfnr}{P10 ’14}{116.}{2 BE}
Bei einer anderen Bank legt Lena 1 000 € für 5 Jahre an; der Zinssatz liegt bei 3 \% pro Jahr. Berechne das Guthaben auf diesem Konto nach den 5 Jahren.
\pfzwfrage{Gib den Wachstumsfaktor für 3 \% Zinsen pro Jahr an.}
\rechenplatz{3}
\fusshilfe{116: 1000 € · 1,03$^{5}$ $\approx$ 1159,27 €}
\end{pfnr}
\pfgruppe{}
\begin{pfnr}{eigene Aufgabe}{117.}{}
Lara bekommt zur Jugendweihe 2\,000\,€ und legt sie zu 3\,\% mit Zinseszins an. Der Führerschein kostet in 4 Jahren etwa 2\,250\,€. Reicht das Geld für den Führerschein?
\par\noindent Antwort: \leerfeld \par
\rechenplatz{2}
\fusshilfe{117: Ja, knapp: nach 4 Jahren $\approx$ 2\,251,02\,€}
\end{pfnr}
\pfpruefstein{P10 ’23}
\begin{pfnr}{}{}{}
Im Jahr 2020 gab es auf der Erde rund 7,8 Milliarden Menschen. Die Tabelle teilt sie nach dem Alter auf: Kinder (0 bis 14 Jahre) 1,95 Milliarden, Jugendliche (15 bis 24 Jahre) 1,25 Milliarden, mittleres Alter (25 bis 64 Jahre) 3,9 Milliarden, höheres Alter (ab 65 Jahre) 0,7 Milliarden.
\end{pfnr}
\begin{pfnr}{}{a)}{2 BE}
Berechne den Anteil der Jugendlichen an der Weltbevölkerung im Jahr 2020 in Prozent.
\rechenplatz{3}
\end{pfnr}
\begin{pfnr}{}{b)}{2 BE}
Eine der beiden folgenden Aussagen widerspricht der Tabelle: Kreuze die falsche Aussage an. Formuliere eine berichtigte Aussage, die zur Tabelle passt.
\pfkreuzzeile{„2020 war jeder zweite Mensch auf der Erde im mittleren Alter.“}\pfkreuzzeile{„2020 war ein Drittel aller Menschen im Kindesalter.“}
\rechenplatz{3}
\end{pfnr}
\end{document}